\runninghead{The Interpretations Compared}
\section{The Four Interpretations Compared}
\label{sec:comparison}

The first two interpretations disagree about the Gospel's signs. For Gregory
and Cyril of Alexandria they are coming portents in the heavens, blows on an
aging world, and the disciple rejoices because what he loves is not ending.
For Ambrose and Augustine they are figures of the Church in persecution, in
which many fall away, and Christ's coming is already under way in the Church
(Augustine) and, Ambrose suggests with a ``perhaps,'' in each one who
receives him. The two interpretations direct the hearer's attention to different
things: Gregory to the ruin of the world, which he must not love; Ambrose and
Augustine to the state of faith, in himself and in the Church. They also
treat the news differently. Gregory takes the disasters of his time as
heralds of the judge, and Augustine warns that to date the end from such
events is a dangerous error. Gregory's own distinction between the herald
and the judge narrows that disagreement without removing it.

The third interpretation begins from the Epistle and does not depend on what
the signs are. Its urgency comes from the sleeper's condition and from the
nearness of salvation, which Aquinas finds in time, in the growth of charity,
and, as the Church reads Paul in Advent, in the first coming. It adds a point
the first two do not make: a man may know that it is the hour and still say
``presently,'' as Augustine did, and the Collect asks Christ who comes to set
us free.

The fourth begins from the Psalm 84 chants. It shares with the third a view
toward Christmas. The third is concerned with the hearer, who must wake; the
fourth with the gift itself, Christ asked for although he has already been
given, given once as man and to be given as God. It alone sets the two
readings of the Communion's earth against each other, the Virgin who bears
Christ and the believer who bears good works; the second and third
interpretations read the fruit as good works without the dispute.

\begin{comparisontable}{Interpretation}{The Gospel's signs}{What is near}{What the hearer is told}
Lift up your heads (Gregory; Cyril of Alexandria) &
Real portents on an aging world; creation begins to change &
Redemption: the body's incorruption and creation renewed &
Do not love what is ending; fear the judge now, and rejoice at his coming\\
The lights in the Church (Augustine; Ambrose) &
Figures of the Church in persecution; each Father assigns sun and stars his
own way &
Christ's coming, already under way in the Church; the hour hidden &
Keep faith, watch daily, and do not try to date the end\\
The hour to rise (Chrysostom; Aquinas) &
Not its subject; the Epistle's night and day &
Salvation, near in time, by charity, and in the Church's Advent memory of
the first coming &
Rise from sin and sloth, put on Christ, and ask Christ to set you free\\
Shew us thy mercy (Augustine; Bellarmine) &
The coming in majesty as the second horizon of the Mass &
Christ, given as man and to be given as God &
Ask for Christ still; bear the fruit God's goodness gives\\
\end{comparisontable}
